
Spurious correlations cause deep neural networks to rely on dataset biases rather than causal features. On the Waterbirds benchmark, models exploit the 95\% correlation between background type (water/land) and bird label, achieving approximately 95\% average accuracy while failing on the worst group (waterbirds on land backgrounds) with substantially lower accuracy. Mitigating spurious correlation reliance without group annotations remains an open problem.

Existing methods such as Group DRO \citep{sagawa2020distributionally} require group annotations during training. Two-stage approaches including Just Train Twice \citep{liu2021just} and Learning from Failure \citep{nam2020learning} exploit early-training signals to identify minority groups without annotations but introduce computational overhead and hyperparameter sensitivity.

This work tests a hypothesis derived from simplicity bias literature: that spurious features, being simpler, emerge not only earlier but also more uniformly across training samples than core features. Specifically, the coefficient of variation (CV) of probe accuracy improvement rates across random sample subsets should be lower for spurious features than for core features. If validated, this signal could enable single-run, annotation-free spurious feature detection.

The hypothesis was tested by training logistic regression probes on frozen CLIP ViT-B/16 features extracted from Waterbirds images. Regularization strength (C) was swept as a proxy for training progression. CV was computed across five random 20\% subsets of training data.

The main finding is negative: CV-based detection achieves AUC = 0.0 on Waterbirds. Both feature types show nearly identical CV values (~0.04), with the direction marginally reversed from the hypothesis (CV of spurious features slightly higher than CV of core features). The root cause is feature saturation in pretrained models---CLIP has already learned both background and bird type concepts, eliminating any emergence dynamics that probing could detect.

\subsubsection{Contributions}

\begin{enumerate}
\item A negative result with clear attribution: CV-based spurious detection fails on frozen pretrained features, achieving AUC = 0.0 on Waterbirds with CLIP ViT-B/16.
\end{enumerate}

\begin{enumerate}
\item Mechanistic explanation: Feature saturation in pretrained models eliminates emergence dynamics; probes converge immediately regardless of regularization strength.
\end{enumerate}

\begin{enumerate}
\item Methodological clarification: Emergence-based detection requires training-time measurement, not post-hoc probing on frozen representations.
\end{enumerate}

